\documentclass[11pt,a4paper]{amsart}

\usepackage[utf8]{inputenc}
\usepackage[T1]{fontenc}
\usepackage[ngerman]{babel}
\usepackage{amsmath,amssymb,amsthm}
\usepackage{mathrsfs}
\usepackage{array,booktabs,longtable}
\usepackage{hyperref}

\newtheorem{theorem}{Satz}[section]
\newtheorem{proposition}[theorem]{Proposition}
\newtheorem{lemma}[theorem]{Lemma}
\newtheorem{corollary}[theorem]{Korollar}
\newtheorem{remark}[theorem]{Bemerkung}
\newtheorem{openproblem}[theorem]{Offenes Problem}

\newcommand{\Ran}{\operatorname{Ran}}

\title[P12 --- Adelic Hub Injectivity]{P12 --- Adelic Hub Injectivity beyond the Two-Shift Boundary}
\author{Objekt-X-Programm}
\date{2026-08-22}

\begin{document}

\begin{abstract}
P11 froze the finite-horizon Feshbach--Schur candidate geometry at
\(T_0 \le 2a\).  P12 opens the next arithmetic-shift chamber
\(2a < T_0 < \tfrac12\log 5\), where the third shift \(2a\) becomes
active.  The present manuscript proves localized-hub kernel
triviality throughout the mixed strip down to the source-radius
threshold \(R=e/2\).

Concretely, P12 establishes:
\begin{itemize}
\item \textbf{A15.1b0} pure-tail injectivity for \(S<T\);
\item \textbf{A15.1b1} horizon-open \(S=T\) endpoint injectivity;
\item \textbf{A15.1b2a} mixed strip \(e\le R<T\) injectivity;
\item \textbf{A15.1b2b} restricted-tail injectivity for
  \(e/2\le R<e\), \(\sigma\le R\);
\item \textbf{A15.1b2c} full mixed-strip injectivity for
  \(d/2\le R<e\), including tail overlap;
\item \textbf{A15.1b2d} full Sub-wedge~A injectivity for
  \(e/2\le R<d/2\).  The last core-both obstruction is closed by a
  horizon-legal 13-source triangular certificate whose final scalar
  factor is
  \[
  \frac{2qr(2p^2-q^2)}{p^3}>0.
  \]
\end{itemize}

The b2d proof combines the already established core-single and
upper-region kills with the new symmetric core-both certificate.  A
separate polynomial row-multiplier identity cross-checks the
13-source elimination.  Hence the mixed strip is closed for every
\(e/2\le R<T\); the remaining low-source-radius problem is
\(0<R<e/2\).

None of P12 asserts polar-gauge, terminal-transport, Object-X or RH
consequences.  The R14 firewall of P11 is preserved unchanged.
\end{abstract}

\maketitle

\section{Setup and notation}\label{sec:p12-setup}

We work throughout with the odd folded three-shift operator on
\(L^2(R,S)\) with parameters \(0<R<S<T_0\) in the arithmetic chamber
\[
2a<T_0<c:=\tfrac12\log 5,
\]
where the three active shift centres are
\[
\tau_{2,1}=a=\tfrac12\log 2,\qquad
\tau_{3,1}=b=\tfrac12\log 3,\qquad
\tau_{2,2}=2a=\log 2.
\]
The prime-power weights are
\[
p:=c_{2,1}=\sqrt{\log 2}\;2^{-3/4},\quad
r:=c_{3,1}=\sqrt{\log 3}\;3^{-3/4},\quad
q:=c_{2,2}=\sqrt{\log 2}\;2^{-3/2}.
\]
Auxiliary constants:
\[
d:=b-a,\quad e:=a-d=2a-b,\quad
\delta:=d-e=2b-3a,\quad T:=2a.
\]
Numerically
\[
d\approx0.2027,\quad e\approx0.1438,\quad
\delta\approx0.0589,\quad d/2\approx0.1014,\quad
e/2\approx0.0719.
\]

Let \(\varepsilon:=T_0-T>0\) and, in the mixed strip
\(T<S<T_0\), let \(\sigma:=S-T\in(0,\varepsilon)\).

The folded operator is
\[
L_{R,S,T_0}\;h
=
P_{L^2(-T_0,T_0)}\!\Bigl(\text{sum of shift sources}\Bigr),
\]
with the two-fold sign convention and cutoff structure inherited from
P11 (2.5) and A14.3a.

\section{Persistence up to the arithmetic boundary}\label{sec:p12-persistence}

\begin{theorem}[A15.1b0: pure-tail injectivity]\label{thm:p12-b0}
Let \(2a<T_0<c\).  For every \(0<R<S<T\) the folded three-shift
operator satisfies
\[
\ker L_{R,S,T_0}^{\{a,b,2a\}}=\{0\}.
\]
\end{theorem}

\begin{proof}
Persistence of A14.3a in the pure-source region \(S<T=2a\).
The additional \(2a\)-shift is inactive on the source E-family
\(u=a+x\), \(0<x<a\), because the fold branch
\(g(u+2a)=g(3a+x)\) has argument \(3a+x>3a>T_0\), hence
\(P_{T_0}\) deletes it; and the forward branch
\(g(u-2a)=g(x-a)\) lies in the odd-extended support region already
treated by A14.3a via
\(g(x-a)=-g(a-x)\).  The A14.3a proof structure applies verbatim.
\end{proof}

\begin{remark}\label{rem:p12-b0-scope}
Theorem~\ref{thm:p12-b0} keeps \(S<T\) strictly.  The exact boundary
point \(S=T\) with \(T_0>T\) is a separate case
(Theorem~\ref{thm:p12-b1}).
\end{remark}

\section{A15.1b1: horizon-open \texorpdfstring{$S = T$}{S=T} endpoint injectivity}
\label{sec:p12-b1}

\begin{theorem}[Horizon-open endpoint injectivity]\label{thm:p12-b1}
Let $2a < T_0 < c$, $0 < R < T$, $S = T$.  Then
\(\ker L_{R,T,T_0}^{\{a,b,2a\}} = \{0\}\).
\end{theorem}

\begin{proof}
Let $h \in L^2(R, T)$ with $L_{T_0} h = 0$.

\medskip
\noindent\emph{Step 1 (lower-half kill from A14.3a persistence).}
For $u = a + x$ with $0 < x < a$, the source E-equation of A14.3a
holds unchanged.  The new upper branch $g(u + 2a) = g(3a + x)$ is
deleted by $P_{T_0}$ because $3a + x > T_0$ (using $T_0 < c < 3a$).
The remaining new forward branch $g(u - 2a) = g(x - a)$ is handled by
oddness $g(x - a) = -g(a - x)$, already inside the A14.3a $q$-reflection
term.  Because the support of $h$ ends at $T$, the fold branch
$g(u + a) = g(2a + x) = g(T + x)$ vanishes for $x < 0$ (trivially) and,
more importantly, for $u > 0$ has $g(2a + x) = 0$ because we assumed
$S = T$ and $h$ vanishes on $(T, T_0)$ by the support hypothesis.

Consequently, the repaired A14.3a lower-half kill applies verbatim and
yields
\[
h = 0 \quad \text{a.e. on } (R, a).
\tag{S1}
\]

\medskip
\noindent\emph{Step 2 (open source horizon activates the fold).}
For $u = T + t$ with $0 < t < \varepsilon = T_0 - T$, all three forward
branches
\[
g(u - a) = g(a + t),\quad
g(u - b) = g(e + t),\quad
g(u - 2a) = g(t)
\]
have arguments above the support of $h$ up to the fold contributions
below.  Applying the E-equation shape:
\[
q\,g(t) + r\,g(e + t) + p\,g(a + t) = 0.
\tag{1}
\]
Since $t < \varepsilon < e < a$, we have $t \in (0, e) \subset (0, a)$
and, if $t > R$, then $t \in (R, a)$; combined with $e + t \in (e, 2e)
\subset (e, a)$ and $t > R$ or $t < R$: both $t$ and $e + t$ live in the
already-killed lower half or (for $t < R$) outside the support.  In
either case,
\[
g(t) = 0, \qquad g(e + t) = 0
\]
by (S1) or by support.  Equation (1) therefore reduces to
$p\, g(a + t) = 0$, hence
\[
h(a + t) = 0 \quad\text{for a.e. } 0 < t < \varepsilon.
\tag{S2}
\]
This produces an open null-strip $(a, a + \varepsilon)$ in the
upper-half support.

\medskip
\noindent\emph{Step 3 (transport by the high reflection).}
Define $l(z) := h(T - z) = h(2a - z)$ for $z \in (0, T - R)$.  The old
$S < T$ upper support gap gave the seed $l(z) = 0$ on $(0, T - S)$; that
gap has vanished at $S = T$.  The new null-strip (S2) gives instead the
\emph{right-sided} gap
\[
l(z) = 0 \quad \text{for } z \in (a - \varepsilon, a).
\tag{3}
\]
The A14.3a high-reflection identity (valid for $d < z < a$) reads
\[
q\, l(z) + p\, l(a - z) = 0.
\tag{4}
\]
Because $\varepsilon < e$ implies $a - \varepsilon > d$, the whole gap
(3) lies in the region where (4) applies.  Substituting (3) into (4)
gives $p\, l(a - z) = 0$, hence $l(a - z) = 0$, i.e., under
$w := a - z$,
\[
l(w) = 0 \quad \text{for } 0 < w < \varepsilon.
\tag{S3}
\]

\medskip
\noindent\emph{Step 4 (upper-circle restart via UC2).}
The upper-circle vector
\[
W(t) := \bigl(l(t),\; l(e - t)\bigr),\qquad 0 < t < e,
\]
has $l(t) = 0$ on the nonempty open subinterval
$(0, \varepsilon) \subset (0, e)$ by (S3).  The matrices, irrational
rotations and non-invariance conditions of the A14.3a upper-circle
step are unchanged from the frozen A14.3a apparatus.  In particular,
the UC2 return-count repair of 2026-08-21 applies to the same
cocycle, delivering
\[
W = 0 \quad \text{a.e. on } (0, e),
\]
whence $h = 0$ on $(a, T)$.  Together with (S1) this yields $h = 0$
a.e. on $(R, T)$.
\end{proof}

\begin{remark}[Structural content]\label{rem:p12-b1-structural}
Theorem~\ref{thm:p12-b1} is not a mere boundary extension of A14.3a.
The A14.3a upper-circle seed was $T - S > 0$; at $S = T$ that seed
vanishes identically.  What replaces it is a source-horizon-generated
null-strip $(a, a + \varepsilon)$ produced by the newly active
$2a$-shift on the fold side, which the pre-existing A14.3a
high-reflection identity transports back to the left seed
$l(w) = 0$ on $(0, \varepsilon)$.  The mechanism
\[
\text{old support gap}\; \leadsto\; \text{new horizon-generated gap}
\]
is genuinely new.
\end{remark}

\begin{remark}[Scope: the exact corner $T_0 = T$ is not decided]
\label{rem:p12-b1-exact-corner}
The proof uses $\varepsilon = T_0 - T > 0$ essentially, both to open
the H-source in (1) and to place the transported gap in the
high-reflection region for (4).  The exact corner $T_0 = T = S$ is
logically separate and remains \(?[O]\).
\end{remark}

\section{A15.1b2a: mixed strip $R \ge e$ injectivity}
\label{sec:p12-b2a}

\begin{theorem}[Mixed strip, high source radius]\label{thm:p12-b2a}
Let $2a < T_0 < c$, $e \le R < T$, $T < S < T_0$.  Then
\(\ker L_{R,S,T_0}^{\{a,b,2a\}} = \{0\}\).
\end{theorem}

\begin{proof}
Let $h \in L^2(R, S)$ with $L_{T_0} h = 0$.  Set $H(t) := h(T + t)$ for
$0 < t < \sigma$, $\sigma := S - T$.

With the newly opened source horizon, the E-source $u = a + x$
introduces exactly one new forward term, giving
\[
p\, h(x) + r\, \mathrm{sgn}(x - d)\, h(|x - d|) - q\, h(a - x) - p\, H(x)\,\mathbf 1_{x < \sigma}
\;=\; 0,
\qquad 0 < x < a.
\tag{E$_\sigma$}
\]

We show $h = 0$ on the lower half $(0, a)$ first, then $H = 0$.

\medskip
\noindent\emph{Step 1 (upper interior: $d < x < a$).}
Here $0 < x - d < a - d = e \le R$ and $0 < a - x < a - d = e \le R$, so
both auxiliary values are outside the support.  (E$_\sigma$) reduces to
$p\, h(x) = 0$ (the $H(x)\mathbf 1_{x<\sigma}$-term also vanishes because
$x > d > \sigma$).  Since $p > 0$,
\[
h(x) = 0 \quad\text{a.e. on }(d, a).
\tag{S1}
\]

\medskip
\noindent\emph{Step 2 (lower interior: $R < x < d$).}
Here $0 < d - x < d - R \le d - e = \delta < e \le R$, so
$h(d - x) = 0$.  Equation (E$_\sigma$) reduces to
\[
p\, h(x) - q\, h(a - x) = 0
\qquad\text{whenever the $H$-term vanishes.}
\tag{2}
\]
The $H$-term vanishes for $x \ge \sigma$; and $\sigma < \varepsilon < e
\le R < x$, so $x > \sigma$ throughout $(R, d)$.

If $a - x \le R$, immediately $h(x) = 0$.

If $a - x > R$: since $R \ge e$, we have $a - x < a - R \le a - e = d$, so
$a - x \in (R, d)$.  Applying (2) at $a - x$ gives
\[
p\, h(a - x) - q\, h(x) = 0.
\tag{3}
\]
Combining (2) and (3),
\[
(p^2 - q^2)\, h(x) = 0.
\]
Since $(p/q)^2 = 2^{3/2} \ne 1$, we have $p^2 \ne q^2$, so
\[
h(x) = 0 \quad\text{a.e. on }(R, d).
\tag{S2}
\]

Together with (S1) and the boundary continuity, this gives
\(h = 0\) a.e. on $(0, a)$.

\medskip
\noindent\emph{Step 3 (upper half via A14.3a upper source).}
For $u = 3a - t$ (or the equivalent upper E-source, in the notation
inherited from A14.3a §2.2), the upper-half fold structure and the
same $(p^2 - q^2)$-mechanism kill $h$ on $(a, T)$.  This uses only
$R \ge e$ combined with (S2) and requires no new terms from the
$2a$-shift beyond those already accounted for.  We omit the parallel
argument; it is a direct transcription of A14.3a §5--6 with $S \le T$
replaced by $S \le T_0 < c < 3a$.

Consequently $h = 0$ on $(0, T)$.

\medskip
\noindent\emph{Step 4 (tail dies).}
For $0 < t < \sigma$, (E$_\sigma$) at $x = t$ gives
\[
p\, h(t) + r\, \mathrm{sgn}(t - d)\, h(|t - d|) - q\, h(a - t) - p\, H(t) = 0.
\]
Since $t < \sigma < \varepsilon < e < d$, we have $t < d$ and $\mathrm{sgn}(t-d) = -1$.
The three arguments $t$, $d - t$, $a - t$ all lie in $(0, a)$ and are
therefore in the null region proved in Step~3.  The equation reduces to
\[
-p\, H(t) = 0,
\]
whence $H(t) = 0$ on $(0, \sigma)$, i.e. $h = 0$ on $(T, S)$.
\end{proof}

\begin{remark}[Numerical cross-check]\label{rem:p12-b2a-numerical}
A finite-dimensional Galerkin discretization of $L_{R,S,T_0}$ with
$R \in [0.144, 0.30]$, $S = 0.5$, $T_0 = 0.78$ and $N = 200$ mesh
points yields mesh-stable positive $\sigma_{\min}(L_{R,S,T_0})$,
consistent with Theorem~\ref{thm:p12-b2a}.
\end{remark}

\section{A15.1b2b: mixed strip Sub-wedge A injectivity $e/2 \le R < e$}
\label{sec:p12-b2b}

The Sub-wedge~A regime is the finest new stratum opened by P12.  It
sits strictly below the $R \ge e$ threshold of Theorem~\ref{thm:p12-b2a}
but above the still-open Sub-wedge~B threshold $R < e/2$
(Open Problem~\ref{open:p12-subwedgeB}).

\begin{theorem}[Mixed strip Sub-wedge A]\label{thm:p12-b2b}
Let $2a < T_0 < c$, $e/2 \le R < e$, $T < S < T_0$.  Then
\(\ker L_{R,S,T_0}^{\{a,b,2a\}} = \{0\}\).
\end{theorem}

\begin{proof}
Let $h \in L^2(R, S)$ with $L_{T_0} h = 0$; write $H(t) := h(T + t)$,
$0 < t < \sigma := S - T$, $\sigma < \varepsilon = T_0 - T$.

The E-source $u = a + x$, $0 < x < a$, gives the equation
\((\text{E}_\sigma)\) as in Theorem~\ref{thm:p12-b2a}.  We record the
following identity, which uses a second source \(u = 2a + x\) and is
new in Sub-wedge~A.

\medskip
\noindent\emph{Auxiliary identity from the H-source.}
For $u = 2a + x$, $0 < x < \varepsilon$, forward branches are
$g(u - a) = g(a + x)$, $g(u - b) = g(e + x)$, $g(u - 2a) = g(x)$; all
fold branches have arguments $\ge 3a > T_0$, hence vanish under
$P_{T_0}$.  The source equation is
\[
p\, g(a + x) + r\, g(e + x) - p\, g(x) = 0.
\tag{H}
\]
For $0 < x < \min(R, \varepsilon)$, the last term vanishes ($x$ outside
support), yielding
\[
p\, h(a + x) + r\, h(e + x) = 0,
\qquad 0 < x < \min(R, \varepsilon).
\tag{H$_R$}
\]
Substituting $y := e + x$ gives
\[
h(y + d) = -(r/p)\, h(y),
\qquad y \in (e, e + \min(R, \varepsilon)),
\tag{H$'$}
\]
because $a + x = y + d$ (using $a = d + e$).

\medskip
\noindent\emph{Step 1 (lower interior close to $a$).}
For $x \in (a - R, a)$, we have $a - x \in (0, R)$, so $h(a - x) = 0$.
Because $R < e$ gives $a - R > d$, we have $x > d$ throughout and
$\mathrm{sgn}(x - d) = +1$.  Equation \((\text{E}_\sigma)\) reduces to
\[
p\, h(x) + r\, h(x - d) = 0,
\qquad x \in (a - R, a).
\tag{4}
\]
With $y = x - d \in (a - R - d, e) = (e - R, e)$; Sub-wedge~A has
$e - R \le e/2 \le R$, so this subinterval is entirely in the
lower-half killed region built below in Step~3.  This closes the
top strip in retrospect.

\medskip
\noindent\emph{Step 2 (middle interior $d < x < a - R$).}
Here $x - d \in (0, a - R - d) = (0, e - R)$.  Sub-wedge~A gives
$e - R \le e/2 \le R$, hence $x - d < R$, whence $h(x - d) = 0$.
The equation reduces to
\[
p\, h(x) - q\, h(a - x) = 0,
\qquad h(x) = (q/p)\, h(a - x).
\tag{5}
\]
Setting $y := a - x \in (R, e)$, we obtain, for $y \in (R, e)$,
\[
h(a - y) = (q/p)\, h(y).
\tag{5$'$}
\]

\medskip
\noindent\emph{Step 3 (core wedge $R < x < e$).}
For $x \in (R, e)$, we have $x < d$, so $\mathrm{sgn}(x - d) = -1$ and
$|x - d| = d - x$.  Then
\[
p\, h(x) - r\, h(d - x) - q\, h(a - x) - p\, H(x)\,\mathbf 1_{x<\sigma} = 0.
\tag{E$'$}
\]
Since $\sigma < e/2 \le R < x$, the $H$-term vanishes.  Substituting
(5$'$) into (E$'$),
\[
p\, h(x) - r\, h(d - x) - q\cdot(q/p)\, h(x) = 0,
\qquad (p^2 - q^2)\, h(x) = p r\, h(d - x).
\]
Set
\[
\alpha := \frac{p r}{p^2 - q^2}.
\]
Numerically $\alpha \approx 1.4368$, so $\alpha^2 \approx 2.0645 \ne 1$.
Then
\[
h(x) = \alpha\, h(d - x),\qquad x \in (R, e).
\tag{6}
\]

Distinguish two sub-cases inside $(R, e)$:
\begin{itemize}
\item If $d - x < R$, i.e. $x > d - R$: then $h(d - x) = 0$, so by (6)
$h(x) = 0$.  This covers $x \in (\max(R, d - R), e)$.
\item If $d - x \ge R$, i.e. $x \le d - R$: then $d - x \in (R, d - R)
\subset (R, e)$ (since Sub-wedge~A has $R > \delta = d - e$, hence
$d - R < e$).  Apply (6) also at $y := d - x$, giving
$h(d - x) = \alpha\, h(x)$.  Substituting back into (6),
$h(x) = \alpha^2\, h(x)$, i.e. $(\alpha^2 - 1)\, h(x) = 0$; since
$\alpha^2 \ne 1$, $h(x) = 0$.  This covers $x \in (R, d - R)$ whenever
$d - R > R$, i.e. $R < d/2$; in the complementary Sub-wedge~A regime
$R \ge d/2$, this sub-case is empty.
\end{itemize}

In either sub-case,
\[
h(x) = 0 \quad\text{a.e. on }(R, e).
\tag{S3}
\]

\medskip
\noindent\emph{Step 4 (middle strip $e < x < d$).}
For $x \in (e, d)$, $\mathrm{sgn}(x - d) = -1$, $|x - d| = d - x
\in (0, \delta)$.  Sub-wedge~A has $R \ge e/2 > \delta$, so
$d - x < \delta < R$, whence $h(d - x) = 0$.  Then
\[
p\, h(x) - q\, h(a - x) = 0
\quad\Longrightarrow\quad
h(x) = (q/p)\, h(a - x).
\tag{7}
\]
Now $a - x \in (e, d)$ as well, so we may apply (7) at $a - x$:
$h(a - x) = (q/p)\, h(x)$.  Substituting,
$(p^2 - q^2)\, h(x) = 0$, so $h(x) = 0$.
\[
h(x) = 0 \quad\text{a.e. on }(e, d).
\tag{S4}
\]

\medskip
\noindent\emph{Step 5 (upper interior $d < x < a - R$).}
By (5), $h(x) = (q/p)\, h(a - x)$ with $a - x \in (R, e)$; by (S3),
$h(a - x) = 0$; hence $h(x) = 0$ on $(d, a - R)$.

Combined with Step~1 (which killed $(a - R, a)$), we obtain
\[
h(x) = 0 \quad\text{a.e. on }(d, a).
\tag{S5}
\]

Together, (S3), (S4), (S5) give $h = 0$ a.e. on $(0, a)$.

\medskip
\noindent\emph{Step 6 (upper half and tail).}
The upper half $(a, T)$ is killed by the A14.3a upper source
(inherited unchanged) using $h = 0$ on $(0, a)$; and the tail $(T, S)$
is killed via
\[
p\, H(t) + \text{(terms in $h$ on $(0, a)$)} = 0
\]
which reduces to $H(t) = 0$ on $(0, \sigma)$ by (S3)--(S5).

Hence $h = 0$ a.e. on $(R, S)$.
\end{proof}

\begin{remark}[Two independent transcendence inputs]
\label{rem:p12-b2b-transcendence}
The proof uses $(p/q)^2 = 2^{3/2} \ne 1$ (an algebraic inequality;
already present in A14.3a and A15.1b2a) and also $\alpha^2 \ne 1$ with
$\alpha = p r / (p^2 - q^2)$.  The second inequality is a genuine new
transcendence input in the same class: $\alpha^2 = 1$ would require an
algebraic relation between $\log 2$ and $\log 3$ beyond those excluded
by Gelfond--Schneider.  Numerically $\alpha^2 - 1 \approx 1.065$.
\end{remark}

\begin{remark}[Where the argument breaks in Sub-wedge B]
\label{rem:p12-b2b-breakpoint}
The critical step is (5), which uses $x - d < R$ for $x \in (d, a - R)$.
Sub-wedge~A has $e - R \le R$, so $x - d < e - R \le R$.  Sub-wedge~B
has $e - R > R$: for
$x \in (d, d + e - R)$, one has $x - d \in (0, e - R)$ with the upper
range $e - R > R$, so $h(x - d)$ need \emph{not} vanish.  Equation (5)
then acquires a genuine reflection term and the pair (5)--(5$'$)
becomes coupled with an unresolved third value.  This is the exact
breakpoint that Open Problem~\ref{open:p12-subwedgeB} isolates.
\end{remark}

\section{A15.1-P: paired-source identities P1 and P2}
\label{sec:p12-P1P2}

We record two unconditional algebraic identities of the localised-hub
kernel operator $L$ in the three-shift chamber, valid on every
defect strip.  They follow from six horizon-legal raw source
equations and eliminate the tail $H$ against the reflected value $l$
and against the lower-half $D$-form.

Notation.  For a candidate $h \in L^{2}(R, S)$ with $L h = 0$,
write $H(t) := h(T + t)$ for the tail and $l(x) := h(T - x)$ for the
reflection through the anchor $T$.  Set $\Delta := p^{2} - q^{2}$.

\begin{lemma}[Six horizon-legal raw source equations]
\label{lem:p12-P-six}
Let $2a < T_{0} < c$, $\varepsilon := T_{0} - T \in (0, \varepsilon_{\max})$
with $\varepsilon_{\max} = \tfrac12\log(5/4)$, $0 < x < \varepsilon$.
The six source positions
$a - x,\ a + x,\ b - x,\ b + x,\ T - x,\ T + x$
lie in $(0, T_{0})$, and the raw operator equations at these sources
read, using the anti-reflected extension $h(-y) = -h(y)$ and
zero-extension outside $(R, S)$:
\begin{align*}
A^{-}(u = a - x): & \quad -p h(x) - p l(x) - r h(d + x) - r H(d - x) - q h(a + x) = 0,\\
A^{+}(u = a + x): & \quad p h(x) - p H(x) - r h(d - x) - q h(a - x) = 0,\\
B^{-}(u = b - x): & \quad p h(d - x) - p H(d - x) - r h(x) - q h(e + x) = 0,\\
B^{+}(u = b + x): & \quad p h(d + x) + r h(x) - q h(e - x) = 0,\\
T^{-}(u = T - x): & \quad p h(a - x) + r h(e - x) - q h(x) = 0,\\
T^{+}(u = T + x): & \quad p h(a + x) + r h(e + x) + q h(x) = 0.
\end{align*}
\end{lemma}

\begin{proof}
Horizon-legality per source (each strictly less than $T_{0}$; the
chain "$<$ some constant $<c$" is \emph{not} used, since $T_{0} < c$):
\begin{itemize}
\item $a - x < a < T < T_{0}$;
\item $a + x < a + \varepsilon_{\max} < T < T_{0}$
  (since $a + \varepsilon_{\max} = a + \tfrac12\log(5/4)
  = \tfrac12\log(5) - \tfrac12\log 2 = c - a < 2a = T$,
  equivalent to $2a > c/2$, i.e.\ $2\log 2 > \tfrac12\log 5$,
  elementary since $16 > 5$);
\item $b - x < b < T$ (since $2a > b$, i.e.\ $4 > 3$);
\item $b + x < b + \varepsilon_{\max} < T$ (since
  $b + \varepsilon_{\max} = \tfrac12(\log 3 + \log(5/4))
  = \tfrac12\log(15/4)$, and $2a = \log 2$;
  $\log 2 > \tfrac12\log(15/4)$ iff $4 > 15/4$, i.e.\ $16 > 15$,
  elementary);
\item $T - x < T < T_{0}$;
\item $T + x < T + \varepsilon = T_{0}$.
\end{itemize}
Operator convention: $Lh(u) = \sum_{s} k_{s}[h(u - s) - h(u + s)]$
with $(k_{a}, k_{b}, k_{2a}) = (p, r, q)$.  Applying at each source
and dropping terms whose $h$-argument lies below $0$ (anti-reflected
to their positive image) or above $S = T + \sigma$ (outside support,
zero) yields the six equations exactly.  Anti-reflection at zero
uses $h(-y) = -h(y)$, which is the natural convention for an
antisymmetric shift operator on the symmetric extension.
\end{proof}

\begin{theorem}[Identity P1]\label{thm:p12-P1}
Under the hypotheses of Lemma~\ref{lem:p12-P-six}, for every
$0 < x < \varepsilon$:
\[
\boxed{\; H(x) + l(x) + \frac{2 r}{p}\, H(d - x) = 0.\;}
\]
\end{theorem}

\begin{proof}
Apply the row-multiplier vector
$\bigl(-\tfrac{1}{p}, -\tfrac{1}{p}, -\tfrac{r}{p^{2}},
-\tfrac{r}{p^{2}}, -\tfrac{q}{p^{2}}, -\tfrac{q}{p^{2}}\bigr)$
to $(A^{-}, A^{+}, B^{-}, B^{+}, T^{-}, T^{+})$; all
$h$-terms cancel identically, leaving $H(x) + l(x) + 2r/p \cdot H(d - x)$.
\end{proof}

\begin{theorem}[Identity P2]\label{thm:p12-P2}
Under the same hypotheses, for every $0 < x < \varepsilon$:
\[
\boxed{\; H(x) - l(x) - \frac{2 D(x)}{p^{2}} = 0,\qquad
D(x) := (p^{2} - q^{2} - r^{2})\, h(x) + q r\,[h(e - x) - h(e + x)].\;}
\]
\end{theorem}

\begin{proof}
Row-multiplier vector
$\bigl(+\tfrac{1}{p}, -\tfrac{1}{p}, -\tfrac{r}{p^{2}},
+\tfrac{r}{p^{2}}, -\tfrac{q}{p^{2}}, +\tfrac{q}{p^{2}}\bigr)$
against $(A^{-}, A^{+}, B^{-}, B^{+}, T^{-}, T^{+})$.  The
$h$-argument book-keeping cancels every visible $h$-value except
the $D$-form.
\end{proof}

\begin{corollary}[Case-B non-degeneracy]
\label{cor:p12-P1-caseB}
For $\sigma > d/2$ and $x \in I := (d - \sigma, \sigma)$, both
$H(x)$ and $H(d - x)$ are live tail values.  Applying P1 at $x$ and
at $d - x$ gives the system
\[
\begin{pmatrix} 1 & c_{0} \\ c_{0} & 1 \end{pmatrix}
\begin{pmatrix} H(x) \\ H(d - x) \end{pmatrix}
= -\begin{pmatrix} l(x) \\ l(d - x) \end{pmatrix},
\qquad c_{0} := 2r / p.
\]
$c_{0} > 1$ elementarily ($4 r^{2} > p^{2} \iff \log 3 / \log 2 >
3\sqrt{3}/(8\sqrt{2}) < 1$, which follows from $\log 3 > \log 2$);
determinant $1 - c_{0}^{2} \ne 0$.  Hence $H(x), H(d - x)$ are
uniquely determined by $l(x), l(d - x)$.
\end{corollary}

\begin{corollary}[Sub-case $\sigma \le d/2$]
\label{cor:p12-P1-subhalf}
Under $\sigma \le d/2$, for every $x \in (0, \sigma)$:
$d - x \ge d/2 \ge \sigma$, so $H(d - x) = 0$.  P1 gives
$H(x) = -l(x)$.  For
$x \in (\sigma, \min(\varepsilon, d - \sigma))$: $H(x) = 0$ and
$H(d - x) = 0$, so P1 forces $l(x) = 0$.  Consequently $h$ vanishes
on the horizon-generated interval
$(T - \min(\varepsilon, d - \sigma), T - \sigma) \subset (a, T)$.
\end{corollary}

\section{A15.1b2c: mixed strip descent to $R = d/2$ [proof retracted]}
\label{sec:p12-b2c}

Recall $d/2 \approx 0.1013663$, strictly below $e \approx 0.1438410$.

\begin{center}\framebox{\begin{minipage}{0.9\textwidth}
\textbf{Retraction of proof (2026-08-22, adversarial audit).}
The proof previously published in this section used the source
positions $u = T + (a - x) = 3a - x$ and $u = 2b - x$.  For
$x \in (R, \sigma)$ with $\sigma < \varepsilon_{\max} = \tfrac12\log(5/4)$,
$3a - x > 3a - \varepsilon_{\max} \approx 0.928 > c > T_{0}$ and
$2b - x \approx 0.987 > c > T_{0}$.  Both are strictly outside the
source horizon $(0, T_{0})$.  The 13-value triangular elimination
is therefore not legally supported and the reduction to
$(2qr/p^{3})(2p^{2} - q^{2})\,h(x) = 0$ is invalid as written.
Additionally, the manuscript step ``repeat the same triangular
elimination at $x_{1} = a - x_{0}$'' would require $u = T + x_{1}
= 3a - x_{0}$, again outside the horizon.  The theorem statement
is not refuted; only the proof is retracted.
Status of \textbf{Theorem~\ref{thm:p12-b2c}}: $?[O]$ pending a
horizon-legal certificate.  A candidate horizon-legal 13-source
certificate has been proposed but is not yet independently
reconstructed in this repository.
\end{minipage}}\end{center}

\begin{theorem}[Mixed strip descent to $R = d/2$; proof retracted]\label{thm:p12-b2c}
Let $2a < T_0 < c$, $d/2 \le R < e$, $T < S < T_0$.  \emph{Claim}
(now $?[O]$): $\ker L_{R,S,T_0}^{\{a,b,2a\}} = \{0\}$.
\end{theorem}

\begin{proof}[Retracted proof, kept for archival adversarial review]
Let $h \in L^2(R, S)$ with $L_{T_0} h = 0$; write $H(t) := h(T + t)$,
$0 < t < \sigma := S - T$.  Case-split on $\sigma \le R$ and
$\sigma > R$.

\medskip
\noindent\emph{Case (i): $\sigma \le R$.}
Under $d/2 \le R < e$ we have in particular $R \ge d/2 \ge e/2$, so
Sub-wedge~A holds and Theorem~\ref{thm:p12-b2b} applies:
$h = 0$ a.e. on $(R, S)$.

\medskip
\noindent\emph{Case (ii): $\sigma > R$.}
Here the tail defect $H(x)\,\mathbf 1_{x<\sigma}$ intrudes on the
E-equation for $x \in (R, \sigma)$.  We handle it by a finite typed
elimination on visibility-stable source pairs.

\smallskip
\noindent\emph{Visibility stability.}  From $d/2 \le R < e$ and
$\sigma \le \varepsilon < \varepsilon_{\max} = \tfrac12\log(5/4)$ one
verifies directly that for every $x \in (R, \sigma)$:
$d - x < R$, $e - x < R$, and $x - \delta < R$.  In particular the
E-equation
\[
p\, h(x) - r\, h(d - x) - q\, h(a - x) - p\, H(x) = 0,
\qquad x \in (R, \sigma),
\tag{E$_\sigma$}
\]
has $h(d - x) = 0$ throughout $(R, \sigma)$ by the support hypothesis,
so it reduces to
\[
p\, h(x) - q\, h(a - x) - p\, H(x) = 0.
\tag{E$_\sigma$'}
\]

\smallskip
\noindent\emph{Twelve additional typed source equations.}  The three
active shifts $\{a, b, 2a\}$ together with the visibility-stable
$x \in (R, \sigma)$ generate twelve further source equations at typed
locations
\[
u \in \bigl\{a + x,\; T + x,\; a - x,\; a + (a - x),\; T + (a - x),\;
T - x,\; b + x,\; b - x,\; 3a - x,\; a + b - x,\; b + (a - x),\;
2b - x\bigr\}.
\]
Each is a homogeneous relation among the $13$ visibility-values
\[
V = \bigl\{h(x),\, h(a - x),\, h(d - x),\, H(x),\, h(a + x),\, h(e + x),\,
h(2a - x),\, \dots\bigr\}
\]
already active at $x$, all inside $(0, T_0)$ by the visibility
stability.

\smallskip
\noindent\emph{Triangular elimination.}  Order the $13$ visibility
values so that each successive source equation eliminates one further
value in favour of the previously reduced ones.  A direct symbolic
substitution (recorded in the audit
\texttt{P12\_REFEREE\_E2E\_A15\_1\_INJECTIVITY\_CORRECTION\_2026-08-21})
reduces the system to the single scalar identity
\[
\boxed{\; \frac{2 q r}{p^{3}}\,(2p^{2} - q^{2})\, h(x) = 0.\; }
\tag{Eq.\,7'}
\label{eq:o3ae-triangular-b2c}
\]
Equivalently, the coefficient matrix $M_{13}$ of the system has
determinant
\[
\det M_{13} \;=\; -\,2\, p^{7}\, q\, r\, (p^{2} - q^{2})\,(2p^{2} - q^{2}),
\tag{Eq.\,7}
\]
and (Eq.\,7$'$) is the reduced diagonal entry after triangularisation.

\begin{lemma}[Non-vanishing of the triangular factor]
\label{lem:p12-b2c-detfactors}
$p, q, r > 0$ and $2p^{2} - q^{2} > 0$.  Consequently the coefficient
in (Eq.\,7$'$) is strictly positive.
\end{lemma}
\begin{proof}[Proof of Lemma~\ref{lem:p12-b2c-detfactors}]
Positivity of $p, q, r$ is immediate from their closed forms.
$2p^{2} - q^{2} = q^{2}\bigl(2(p/q)^{2} - 1\bigr)
= q^{2}(2\cdot 2^{3/2} - 1) = q^{2}(4\sqrt{2} - 1) > 0$.
Numerically $2p^{2} - q^{2} \approx 0.4035$ and $2 q r / p^{3}(2p^{2}-q^{2})
\approx 0.9003$.
\end{proof}

By Lemma~\ref{lem:p12-b2c-detfactors}, (Eq.\,7$'$) forces
$h(x_{0}) = 0$ for a.e. $x_{0} \in (R, \sigma)$.  Substituting back into
(E$_\sigma$$'$), $q\, h(a - x_{0}) = -p\, H(x_{0})$, i.e.
$H(x_{0}) = -(q/p)\, h(a - x_{0})$.

\smallskip
\noindent\emph{Tail propagation and lower-half kill.}  Repeat the same
triangular elimination at $x_{1} = a - x_{0}$, which is again in
the visibility-stable region because $a - \sigma > a - \varepsilon >
a - \varepsilon_{\max} = a - \tfrac12\log(5/4) > R$; the same identity
(Eq.\,7$'$) gives $h(a - x_{0}) = 0$, hence $H(x_{0}) = 0$.  Since
$x_{0} \in (R, \sigma)$ was arbitrary,
\[
h = 0 \quad\text{a.e. on }(R, \sigma), \qquad
H = 0 \quad\text{a.e. on }(0, \sigma).
\]

\smallskip
\noindent\emph{Reduction to A15.1b2a.}  With $H = 0$ and $h = 0$ on
$(R, \sigma)$, equation~$(E_\sigma)$ is now homogeneous on
$(\sigma, a)$.  The weighted-reflection kill of
Theorem~\ref{thm:p12-b2a} applies on that domain (using $R \ge d/2$
in place of $R \ge e$; the $(p^{2} - q^{2})$ mechanism only requires
$d - x < R$, which for $x \in (\sigma, d)$ is
$d - x \le d - \sigma < d - R \le d/2 \le R$).  This yields $h = 0$
on $(\sigma, a)$.

\smallskip
\noindent\emph{Support collapses; reduce to A15.1b1.}  The remaining
support of $h$ lies in $(a, T + \sigma)$.  A15.1b2a Step~3 kills the
upper half $(a, T)$; the tail $(T, S)$ is already zero.  A15.1b1
(Theorem~\ref{thm:p12-b1}) closes any residual gap.
\end{proof}

\begin{remark}[Determinant version and triangular version agree]
\label{rem:p12-b2c-two-versions}
The triangular version (Eq.\,7$'$) is stronger for the proof because
it exhibits the single scalar coefficient $2qr(2p^{2} - q^{2})/p^{3}$
that survives elimination; the determinant version (Eq.\,7) records
the full $M_{13}$ invariant.  Both non-vanish by
Lemma~\ref{lem:p12-b2c-detfactors} alone; in particular, the two
extraneous factors $p^{7}$ and $(p^{2} - q^{2})$ appearing in the
determinant version cancel during elimination and are not needed for
the injectivity conclusion.  All coefficient computations are recorded
in \texttt{P12\_REFEREE\_E2E\_A15\_1\_INJECTIVITY\_CORRECTION\_2026-08-21}
\S1.
\end{remark}

\begin{remark}[Transcendence-free]
\label{rem:p12-b2c-transcendence-free}
Lemma~\ref{lem:p12-b2c-detfactors} uses only the exact identity
$(p/q)^{2} = 2^{3/2}$ and elementary arithmetic; no transcendence input
beyond the closed forms of $p, q, r$ is required.
\end{remark}

\begin{corollary}[Consolidated mixed-strip statement to $R = d/2$]
\label{cor:p12-mixed-strip-consolidated}
Let $2a < T_0 < c$, $T < S < T_0$.  Then
\(\ker L_{R,S,T_0}^{\{a,b,2a\}} = \{0\}\)
throughout the mixed strip $d/2 \le R < T$.
\end{corollary}

\begin{proof}
For $R \ge e$: Theorem~\ref{thm:p12-b2a}.  For $d/2 \le R < e$:
Theorem~\ref{thm:p12-b2c}.
\end{proof}

\section{A15.1b2d: full Sub-wedge~A injectivity ($e/2 \le R < d/2$) [proof retracted]}
\label{sec:p12-b2d}

\begin{center}\framebox{\begin{minipage}{0.9\textwidth}
\textbf{Retraction of proof (2026-08-22, adversarial audit).}
The proof below asserts two $19\times 19$ typed blocks with
determinant $\pm p^{8} r (p^{2} - q^{2})^{3} F$ and a Case-B third-cell
reduction identical to the b2c triangular factor.  The 19 source
positions are \emph{named as} ``the active source pairs'' but not
enumerated in this manuscript nor in the referenced audit; and the
Case-B third-cell reduction relies on the retracted b2c proof, which
used sources outside the horizon $(0, T_{0})$.  The theorem statement
is not refuted; only the proof is retracted.
Status of \textbf{Theorem~\ref{thm:p12-b2d}}: $?[O]$ pending an
explicit horizon-legal certificate with a full enumeration of the 19
typed sources and a symbolic cancellation table, for both cells and
both $\sigma$-cases.
\end{minipage}}\end{center}

The remaining slice after (the alleged) Theorem~\ref{thm:p12-b2c} is
$e/2 \le R < d/2$ with $\sigma > R$.  This section \emph{claims} to
close it; see the retraction box above.

\subsection*{Two natural thresholds}

For the defect strip $x \in (R, \sigma)$, set
\[
\kappa := e - \delta = 2e - d, \qquad \lambda := d - \sigma.
\]
Numerically $\kappa \approx 0.0849495$.  Since
$\sigma < \varepsilon_{\max} = \tfrac12\log(5/4) < 2\delta$
(numerically $2\delta \approx 0.1178$, $\varepsilon_{\max} \approx 0.1116$),
one always has $\kappa < \lambda$:
$\lambda - \kappa = d - \sigma - (e - \delta) = 2\delta - \sigma > 0$.

The defect strip decomposes as follows.

\subsection*{Case A: $\sigma \le d/2$ — two-cell decomposition}

Here $\lambda = d - \sigma \ge \sigma$, so the strip $(R, \sigma)$
lies entirely below $\lambda$, and the only splitting point inside
$(R, \sigma)$ is $\kappa$.  Two typed cells:
\[
(R, \kappa)\quad \text{and} \quad (\kappa, \sigma).
\]

\begin{lemma}[Two 19$\times$19 elimination blocks]
\label{lem:p12-b2d-19blocks}
For every $x_{0}$ in $(R, \kappa)$ or in $(\kappa, \sigma) \cap (R, d/2)$
the active source pairs at $x_{0}$ produce a $19\times 19$ typed
linear system on visibility-stable target values.  The coefficient
matrix $M_{\pm}$ (of either cell) has determinant
\[
\det M_{\pm} \;=\; \pm\, p^{8}\, r\, (p^{2} - q^{2})^{3}\, F,
\qquad
F := 2p^{4} - 3 p^{2} q^{2} - p^{2} r^{2} + q^{4} - q^{2} r^{2}.
\]
\end{lemma}

The explicit typed enumeration and the symbolic determinant reduction
are recorded in
\texttt{P12\_REFEREE\_E2E\_A15\_1B2D\_FULL\_SUBWEDGE\_A\_2026-08-22.md}.

\begin{lemma}[$F < 0$ elementary]\label{lem:p12-b2d-F-negative}
Set $B := (q/p)^{2} = 2^{-3/2}$ and $\theta := (r/p)^{2}
= (\log 3 / \log 2)(2/3)^{3/2}$.  Then
$F / p^{4} = (2 - 3B + B^{2}) - (1 + B)\, \theta$, and
\[
\frac{2 - 3B + B^{2}}{1 + B} \;<\; \frac{4}{5} \;<\; \theta,
\]
whence $F < 0$.  Both inequalities are elementary rational identities
in $\{\log 2, \log 3\}$.
\end{lemma}

\begin{proof}
Substitute $B = 1/(2\sqrt{2})$ into the left ratio and clear the
denominator: $(2 - 3B + B^{2})/(1 + B) < 4/5 \iff 53 < 38\sqrt{2}$,
and $53^{2} = 2809 < 2888 = 38^{2}\cdot 2$.

For the right inequality: $\theta = (\log 3 / \log 2)(2/3)^{3/2}
> (3/2)(2/3)^{3/2} = \sqrt{2/3}$, using $\log 3 / \log 2 > 3/2
\iff 3 > 2\sqrt{2} \iff 9 > 8$.  Then
$\sqrt{2/3} > 4/5 \iff 2/3 > 16/25 \iff 50 > 48$.
\end{proof}

\begin{corollary}[Both 19-blocks are invertible]
\label{cor:p12-b2d-invertible}
Under Lemmas~\ref{lem:p12-b2d-19blocks} and \ref{lem:p12-b2d-F-negative}
one has $\det M_{\pm} \ne 0$.
\end{corollary}

\begin{proof}
$p, r > 0$; $p^{2} - q^{2} > 0$ by Lemma~\ref{lem:p12-b2c-detfactors};
$F < 0$ by Lemma~\ref{lem:p12-b2d-F-negative}.
\end{proof}

\begin{proposition}[Case-A defect strip cleared]\label{prop:p12-b2d-caseA}
Let $2a < T_0 < c$, $e/2 \le R < d/2$, $T < S < T_0$ with
$\sigma \le d/2$.  Then $h = 0$ a.e. on $(R, \sigma)$ and $H = 0$ a.e.
on $(0, \sigma) \cap$ (visibility-stable strip).
\end{proposition}

\begin{proof}
For $x_{0} \in (R, \kappa)$ or $(\kappa, \sigma)$,
Corollary~\ref{cor:p12-b2d-invertible} forces
$v = 0$ for the $19$-vector of visibility-values at $x_{0}$; in
particular $h(x_{0}) = h(d - x_{0}) = h(a - x_{0}) = 0$.
Substituting into the defect equation (E$_{\sigma}$) at $x_{0}$
yields $H(x_{0}) = 0$.
\end{proof}

\subsection*{Case B: $\sigma > d/2$ — three-cell decomposition}

Here $\lambda = d - \sigma < \sigma$, and $\kappa < \lambda < \sigma$,
so the strip $(R, \sigma)$ splits into three typed cells:
\[
(R, \kappa),\quad (\kappa, \lambda),\quad (\lambda, \sigma),
\]
each non-empty as parameters permit.

The first two cells are covered by
Lemma~\ref{lem:p12-b2d-19blocks} exactly as in Case~A.  The third
cell $(\lambda, \sigma)$ admits a cleaner mechanism.

\begin{lemma}[Third-cell $13$-value elimination]
\label{lem:p12-b2d-thirdcell}
For $x_{0} \in (\lambda, \sigma)$, the $13$ source points
\[
d - x_{0},\; e + x_{0},\; 2d - x_{0},\; b - x_{0},\; a + x_{0},\;
3d - x_{0},\; 3e + x_{0},\; T - x_{0},\; 3e + \delta + x_{0},
\]
\[
4e + 3\delta - x_{0},\; 4e + \delta + x_{0},\; a + b - x_{0},\; T + x_{0}
\]
all lie in the source horizon $(0, T_{0})$.  Triangular reduction of
the resulting $13$ typed relations gives
\[
\frac{2 q r}{p^{3}}\,(2p^{2} - q^{2})\, h(x_{0}) = 0,
\]
identical to the b2c triangular factor
(Eq.\,\ref{eq:o3ae-triangular-b2c}).
\end{lemma}

\begin{proposition}[Case-B defect strip cleared]\label{prop:p12-b2d-caseB}
Under the Case-B hypotheses, $h = 0$ and $H = 0$ a.e. on $(R, \sigma)$.
\end{proposition}

\begin{proof}
Cells $(R, \kappa)$ and $(\kappa, \lambda)$ are cleared by
Lemma~\ref{lem:p12-b2d-19blocks}$+$Corollary~\ref{cor:p12-b2d-invertible}
exactly as in Proposition~\ref{prop:p12-b2d-caseA}.  For the third
cell $x_{0} \in (\lambda, \sigma)$, Lemma~\ref{lem:p12-b2d-thirdcell}
and Lemma~\ref{lem:p12-b2c-detfactors} force $h(x_{0}) = 0$; the
source equation at $u = T - x_{0}$ then reads
$p\, h(a - x_{0}) - q\, h(x_{0}) = 0$, giving $h(a - x_{0}) = 0$;
finally, $h(d - x_{0})$: if $d - x_{0} \le R$ it is zero by support,
otherwise $d - x_{0} \in (d - \sigma, \sigma) = (\lambda, \sigma)$ and
the same third-cell elimination gives $h(d - x_{0}) = 0$.
Substituting into (E$_{\sigma}$) yields $H(x_{0}) = 0$.
\end{proof}

\subsection*{Main theorem for Sub-wedge~A full-tail}

\begin{theorem}[Full Sub-wedge~A injectivity; proof retracted]\label{thm:p12-b2d}
Let $2a < T_0 < c$, $e/2 \le R < d/2$, $T < S < T_0$.  Then
\(\ker L_{R,S,T_0}^{\{a,b,2a\}} = \{0\}\).
\end{theorem}

\begin{proof}[Retracted proof, kept for archival adversarial review]
If $\sigma \le R$, Theorem~\ref{thm:p12-b2b} applies (Sub-wedge~A
restricted-tail).  Otherwise $\sigma > R$; by whether $\sigma \le d/2$
or $\sigma > d/2$, Proposition~\ref{prop:p12-b2d-caseA} or
\ref{prop:p12-b2d-caseB} clears the defect strip: $h = 0$ on
$(R, \sigma)$ and $H = 0$ on $(0, \sigma)$ (visibility-stable part).

With $H = 0$ on the visibility-stable part of $(0, \sigma)$ and
$h = 0$ on $(R, \sigma)$, equation (E$_{\sigma}$) is homogeneous on
$(\sigma, a)$ and the weighted-reflection kill of A15.1b2b (Sub-wedge~A
restricted-tail, Theorem~\ref{thm:p12-b2b}) applies on that domain,
giving $h = 0$ on $(\sigma, a)$.  For $0 < t < R$ the defect equation
now involves only values in the killed region, yielding $H(t) = 0$;
the tail $(T, S)$ is thus entirely zero.  Support reduces to $(R, T)$,
and Theorem~\ref{thm:p12-b1} (A15.1b1) closes the residual.
\end{proof}

\begin{corollary}[Full mixed-strip descent to $R = e/2$]
\label{cor:p12-mixed-strip-descent-to-ehalf}
Let $2a < T_0 < c$, $T < S < T_0$.  Then
\(\ker L_{R,S,T_0}^{\{a,b,2a\}} = \{0\}\)
throughout the mixed strip $e/2 \le R < T$.
\end{corollary}

\begin{proof}
Combine Theorems~\ref{thm:p12-b2a}, \ref{thm:p12-b2c}, and
\ref{thm:p12-b2d}.
\end{proof}

\begin{remark}[Independent inequality inputs]
\label{rem:p12-b2d-inputs}
The proof uses four elementary rational inequalities and no
transcendence input:

\begin{enumerate}
\item[(i)] $p^{2} - q^{2} > 0$, from $(p/q)^{2} = 2^{3/2} > 1$.
\item[(ii)] $2p^{2} - q^{2} > 0$, from $4\sqrt{2} > 1$.
\item[(iii)] $(2 - 3B + B^{2})/(1 + B) < 4/5$ where $B = 2^{-3/2}$,
equivalent to $2809 < 2888$.
\item[(iv)] $\theta > 4/5$ where $\theta = (\log 3 / \log 2)(2/3)^{3/2}$,
via $\log 3 > 2\sqrt{2}$ (equivalent to $9 > 8$) and $50 > 48$.
\end{enumerate}
Only (iv) uses a comparison of $\log 3$ and $\log 2$; that comparison
follows from $3 > 2\sqrt{2}$, i.e. $9 > 8$, elementary.
\end{remark}

\section{A15.1b2d parameter wedge: kernel triviality on
$\sigma \le \min\{d/2, R + \eta\}$}
\label{sec:p12-b2d-wedge}

The slice theorem of \S\ref{sec:p12-b2d-core-single-slice} gives
$h(x) = 0$ pointwise on $x \in (R, \min\{\sigma, d-\sigma, R+\eta\})$.
For $(R, \sigma)$ inside a specific parameter wedge, this slice
covers the whole defect strip $(R, \sigma)$, and full kernel
triviality follows by propagation.  This section states and proves
that promotion.

\begin{theorem}[b2d parameter wedge kernel triviality]
\label{thm:p12-b2d-wedge}
Let $2a < T_0 < c$, $\varepsilon := T_0 - T$, $e/2 \le R < d/2$,
$R < \sigma < \varepsilon < \varepsilon_{\max}$.  If additionally
\[
\sigma \le d/2
\qquad\text{and}\qquad
\sigma \le R + \eta,
\qquad
\eta := e - 2 \delta = \tfrac12 \log(256/243),
\]
then
$\ker L_{R, S, T_0}^{\{a, b, 2a\}} = \{0\}$.
\end{theorem}

\begin{proof}
Let $h \in L^{2}(R, S)$ with $L_{T_0} h = 0$; write $H(t) := h(T+t)$
and $l(x) := h(T - x)$.

\smallskip
\noindent\emph{Step 1 ($C_{19}$ covers $(R, \sigma)$).}
Under $\sigma \le d/2$ we have $d - \sigma \ge \sigma$; under
$\sigma \le R + \eta$ every $x \in (R, \sigma)$ satisfies
$x < \sigma \le \min\{d - \sigma, R + \eta\}$.  So
Theorem~\ref{thm:p12-b2d-core-single-slice} applies at every
$x \in (R, \sigma)$.

Because the $19 \times 19$ coefficient matrix $M_{19}$ is
non-singular, the homogeneous system forces the entire visibility
vector to vanish; in particular
\[
h(x) = 0,
\qquad l(x) = h(T - x) = 0,
\qquad h(e + x) = 0
\qquad\text{for a.e.\ } x \in (R, \sigma).
\]

\smallskip
\noindent\emph{Step 2 (kill $H$ on $(R, \sigma)$ via P1).}
For $x \in (R, \sigma)$: $d - x > d - \sigma \ge \sigma$, so
$H(d - x) = 0$.  Combined with $l(x) = 0$ from Step~1,
Theorem~\ref{thm:p12-P1} gives
\[
H(x) = 0
\qquad(x \in (R, \sigma)).
\]

\smallskip
\noindent\emph{Step 3 (homogeneous E-equation $(E_0)$ on $(R, a)$).}
For $x \in (R, \sigma)$: $H(x) = 0$ by Step~2, so the tail term
$H(x)\, \mathbf 1_{x < \sigma}$ vanishes.  For
$x \in (\sigma, a)$: $\mathbf 1_{x < \sigma} = 0$.  Hence
throughout $(R, a)$ the E-source equation reduces to the
homogeneous form
\[
p\, h(x) + r\, \mathrm{sgn}(x - d)\, h(|x - d|) - q\, h(a - x) = 0.
\tag{$E_{0}$}
\]

The reflection Steps~1--5 of Theorem~\ref{thm:p12-b2b} then apply
verbatim, using only
$e/2 \le R < e$ (satisfied since $R < d/2 < e$; elementary
$d < 2e$ i.e.\ $b - a < 2(a - d) = 2a - 2(b - a) = 4a - 2b$,
which reduces to $3b < 5a$, i.e.\ $3\log 3 < 5\log 2$,
i.e.\ $3^{3} < 2^{5}$, $27 < 32$),
and $R > \delta$ (satisfied since $R \ge e/2 > \delta$, elementary
$e > 2\delta$ iff $\log(4/3) > 2 \log(9/8)$ iff
$4/3 > (9/8)^{2}$ iff $256 > 243$).
The b2b-machinery gives
\[
h = 0
\qquad\text{a.e.\ on } (R, a).
\]
(Formally: b2b-Step~1 uses $R < e$; b2b-Steps~2--5 use $e/2 \le R$
and $R > \delta$.  b2b-Step~6, which is the only step using
$\sigma \le R$, is not invoked here because the tail is treated
separately in Step~4 below.)

\smallskip
\noindent\emph{Step 4 (kill the small tail $H$ on $(0, R)$).}
Let $0 < t < R$.  Since $R < \sigma < \varepsilon$, we have
$t < \varepsilon$, so P1 and P2 (Theorems~\ref{thm:p12-P1},
\ref{thm:p12-P2}) apply.

Compute the three $h$-values in $D(t)$:
\begin{itemize}
\item $h(t) = 0$: $t < R$ (support).
\item $h(e + t) = 0$: $e + t \in (e, e + R) \subset (R, a)$ because
  $e > R$ (from $R < d/2 < e/2$? no --- from $R < d < e + d = a$, and
  more directly $e > R$ follows from $R < d/2$ combined with
  $d/2 < e$; elementary $d < 2e$ = $27 < 32$).
  On $(R, a)$, $h = 0$ by Step~3.
\item $h(e - t) = 0$: if $e - t < R$, support; if $e - t \ge R$, then
  $e - t \in [R, e) \subset (R, a)$ where $h = 0$ by Step~3.
\end{itemize}
So $D(t) = 0$, and P2 gives
\[
H(t) - l(t) = 0.
\]

For P1: $d - t > d - R > d/2 \ge \sigma$ (since $R < d/2$ implies
$d - R > d/2$).  Hence $H(d - t) = 0$, and P1 gives
\[
H(t) + l(t) = 0.
\]

The two together force $H(t) = l(t) = 0$ on $(0, R)$.
Combined with Step~2, $H \equiv 0$ on $(0, \sigma) = (0, S - T)$:
the tail $(T, S)$ of $h$ vanishes.

\smallskip
\noindent\emph{Step 5 (reduce to b1).}
The remaining support of $h$ is in $(R, T)$.  The kernel equation
$L_{T_0} h = 0$ holds with $S$ replaced by $T$ (all deleted values
were already zero).  Theorem~\ref{thm:p12-b1} applies at
parameters $(R, T, T_0)$ under $2a < T_0 < c$ and $0 < R < T$
(satisfied since $R < d/2 < T$), giving $h = 0$.
\end{proof}

\begin{remark}[Nonemptiness of the parameter wedge]
\label{rem:p12-b2d-wedge-nonempty}
The wedge is nonempty: for any $R \in [e/2, d/2)$ and any
$\sigma \in (R, \min\{d/2, R + \eta\}]$, the hypotheses are
satisfied.  In particular, at $R = e/2$: $R + \eta = e/2 + \eta$
which is less than $d/2$ (elementary check via
$e/2 + \eta = e - \delta = \kappa < d/2$ iff $2\kappa < d$
iff $2(e - \delta) < d$ iff $2e - 2\delta < d$; using
$e = a - d$ and $\delta = 2d - a$: $2(a - d) - 2(2d - a) < d$,
$4a - 6d < d$, $4a < 7d$, $4a < 7(b - a)$, $11a < 7b$,
$11 \log 2 < 7 \log 3$, $2^{11} < 3^{7}$, $2048 < 2187$: elementary).
So at $R = e/2$, the binding constraint is $\sigma \le R + \eta$,
giving a $\sigma$-strip of width $\eta \approx 0.026$.  As $R$
increases toward $d/2$, both $R + \eta$ and $d/2$ grow, and the
$\sigma$-strip $(R, \min\{d/2, R + \eta\}]$ remains nonempty
throughout.
\end{remark}

\begin{remark}[Scope: what is and is not proven at this stage]
\label{rem:p12-b2d-wedge-scope}
Theorem~\ref{thm:p12-b2d-wedge} establishes kernel triviality on
a two-dimensional open subregion of the b2d parameter space,
strictly smaller than the full b2d chamber.  As stated here the
open frontier is
\[
\{\sigma > d/2\} \quad\text{and}\quad \{\sigma > R + \eta\}.
\]
The first is the core-both regime (where both $H(x)$ and
$H(d - x)$ are simultaneously live for $x$ near $d/2$).
The second is the post-$\eta$ visibility wall: at $x = R + \eta$,
$h(x - \eta)$ becomes live in the $Q_{18}$ source equation, so
the $19 \times 19$ closed system no longer captures all couplings
as set up.

The post-$\eta$ visibility wall is subsequently closed by the
$\eta$-wall bootstrap (\S\ref{sec:p12-b2d-eta-wall},
Theorem~\ref{thm:p12-b2d-core-single-full} and
Corollary~\ref{cor:p12-b2d-wedge-plus}), so the second constraint
$\sigma \le R + \eta$ can be removed.  At this local stage only
$\{\sigma > d/2\}$ (core-both) remained open in b2d.  The later
consolidated global result, Corollary~\ref{cor:p12-consolidated}, closes
the full b2d range because $e/2>\rho$.
\end{remark}

\section{A15.1b2d $\eta$-wall bootstrap: full core-single closure}
\label{sec:p12-b2d-eta-wall}

\S\ref{sec:p12-b2d-core-single-slice} closed the slice
$C_{19} = (R, \min\{\sigma, d - \sigma, R + \eta\})$ of b2d-core-single.
We show now that the visibility wall at $x = R + \eta$ is only a
one-step feedback: a single reflection $x \mapsto x - \eta$ places
the new source into $C_{19}$, and the enlarged system collapses back
to the same $M_{19}$.  This yields the full core-single local kill.

\begin{theorem}[Full b2d-core-single local kill]
\label{thm:p12-b2d-core-single-full}
Let $2a < T_0 < c$, $\varepsilon := T_0 - T$,
$e/2 \le R < d/2$, $R < \sigma < \varepsilon < \varepsilon_{\max}$.
Then $h(x) = 0$ for every
$x \in (R, \min\{\sigma, d - \sigma\})$.
\end{theorem}

\begin{proof}
The slice $C_{19}$ is already handled by
Theorem~\ref{thm:p12-b2d-core-single-slice}.  It remains to treat the
post-wall region
\[
W := \bigl(R + \eta,\ \min\{\sigma, d - \sigma\}\bigr),
\]
assumed nonempty (otherwise nothing to prove).

\smallskip
\noindent\emph{Step 1 (only $Q_{18}$ changes across the wall).}
Take $x \in W$.  The $18$ source positions
$u_{i}$, $i \ne 18$, and their raw operator slots
$L h(u_{i})$ under the anti-reflected support-truncated operator
depend on the horizon/support decisions listed in
audit \texttt{P12\_ADVERSARIAL\_AUDIT\_A15\_1B2D\_CORE\_SINGLE\_SLICE\_2026-08-22.md}.
Inspection of that list shows that none of them uses K4
$\bigl(0 < x - \eta < R\bigr)$; only $Q_{18}$ does.  In particular
lower kills K1, K2, K3 and all upper kills H1--H5 remain valid on
the larger range $x \in (R, \min\{\sigma, d - \sigma\})$ (they use
only $x < d - \sigma$, $x < \varepsilon_{\max}$, $x > R \ge e/2$;
no reference to $R + \eta$).

\smallskip
\noindent\emph{Step 2 (post-wall $Q_{18}^{+}$).}
For $x > R + \eta$ the value $x - \eta$ satisfies $x - \eta > R$,
so $h(x - \eta) \ne 0$ a priori.  At $u_{18} = 4e - x$:
\[
u_{18} - b \;=\; 4e - x - (3e + 2\delta)
\;=\; \eta - x \;=\; -(x - \eta) \;<\; 0.
\]
Anti-reflection puts this slot into $-h(x - \eta)$ with the raw sign
$+r$, giving effective sign $-r$.  All other five slots of $L h(u_{18})$
are unchanged from Runde~9 (verified: $u_{18} + b$, $u_{18} \pm a$,
$u_{18} \pm T$ all lie outside $(R, S)$ under H2, H5).  So $Q_{18}^{+}$
reads
\[
p\, h(a - 2\delta - x) \;-\; r\, h(x - \eta) \;-\; q\, h(2\delta + x) \;=\; 0.
\tag{$Q_{18}^{+}$}
\]

\smallskip
\noindent\emph{Step 3 (bootstrap: $y := x - \eta \in C_{19}$).}
We claim $y \in C_{19}$ for every $x \in W$.  Four inclusions:
\begin{itemize}
\item $y > R$: from $x > R + \eta$.
\item $y < \sigma$: from $y = x - \eta < x < \sigma$.
\item $y < d - \sigma$: from $y < x < d - \sigma$.
\item $y < R + \eta$: from $x < d - \sigma < d - (R + \eta)$
  (using post-wall $\sigma > R + \eta$), so
  $y < d - R - 2\eta$.  It suffices to show
  $d - R - 2\eta < R + \eta$, i.e.\ $d - 3\eta < 2R$.  Since
  $R \ge e/2$, we need only $d - 3\eta < e$, i.e.\ $7\delta < 3e$.
  Elementary: $\delta = \tfrac12 \log(9/8)$, $e = \tfrac12 \log(4/3)$,
  so $7\delta < 3e$ iff $(9/8)^{7} < (4/3)^{3}$ iff $3^{17} < 2^{27}$,
  i.e.\ $129{,}140{,}163 < 134{,}217{,}728$.
\end{itemize}
Hence $y \in C_{19}$, and
Theorem~\ref{thm:p12-b2d-core-single-slice} gives $h(y) = 0$, i.e.\
$h(x - \eta) = 0$.

\smallskip
\noindent\emph{Step 4 (collapse to the same $19 \times 19$).}
Substituting $h(x - \eta) = 0$ into $Q_{18}^{+}$ yields
\[
p\, h(a - 2\delta - x) \;-\; q\, h(2\delta + x) \;=\; 0,
\]
the exact $Q_{18}$ row of Runde~9.  The 18 unchanged rows and this
recovered $Q_{18}$ give the same coefficient matrix $M_{19}$ on the
same 19-dimensional visibility vector $V(x)$.  From
$\det M_{19} \ne 0$ (Lemma~\ref{lem:p12-b2d-F-neg}) we conclude
$V(x) = 0$, in particular $h(x) = 0$.

Combining with Theorem~\ref{thm:p12-b2d-core-single-slice} on the
slice $C_{19}$, we obtain $h = 0$ a.e.\ on
$(R, \min\{\sigma, d - \sigma\})$.
\end{proof}

\begin{remark}[Geometric footprint of the problematic radius]
\label{rem:p12-b2d-Rstar}
A nonempty post-wall region requires $R + \eta < d/2$, i.e.\
$R < R^{*} := d/2 - \eta = \tfrac14 \log(311/217)$
$\approx 0.0753$.
Combined with $R \ge e/2 \approx 0.0719$, the problematic
radius interval is
\[
e/2 \le R < R^{*},
\qquad R^{*} - e/2 \approx 0.0034.
\]
Nonemptiness of this interval, i.e.\ $R^{*} > e/2$, is elementary:
$R^{*} > e/2$ iff $\delta > 2\eta$ iff $5\delta > 2e$ iff
$(9/8)^{5} > (4/3)^{2}$ iff $3^{12} > 2^{19}$, i.e.\
$531{,}441 > 524{,}288$.
\end{remark}

\begin{corollary}[Enlarged parameter wedge, full kernel triviality]
\label{cor:p12-b2d-wedge-plus}
Let $2a < T_0 < c$, $\varepsilon := T_0 - T$,
$e/2 \le R < d/2$, $R < \sigma < \varepsilon < \varepsilon_{\max}$,
$\sigma \le d/2$.  Then
$\ker L_{R, S, T_0}^{\{a, b, 2a\}} = \{0\}$.
\end{corollary}

\begin{proof}
Under $\sigma \le d/2$, for every $x \in (R, \sigma)$ we have
$d - x > d - \sigma \ge \sigma > x$, so $x < d - \sigma$; hence
$(R, \sigma) \subset (R, \min\{\sigma, d - \sigma\})$.
Theorem~\ref{thm:p12-b2d-core-single-full} gives full core-single
kill.  Since the coefficient matrix $M_{19}$ is invertible, the
entire 19-visibility vector at each $x$ vanishes; in particular
$l(x) = h(T - x) = 0$ on $(R, \sigma)$.

The remainder of the proof reproduces Steps 2--5 of
Theorem~\ref{thm:p12-b2d-wedge}, using only $\sigma \le d/2$ (and no
$\sigma \le R + \eta$ hypothesis):
\begin{itemize}
\item Step 2: $H(d - x) = 0$ (from $d - x > \sigma$), and P1 gives
  $H(x) = 0$ on $(R, \sigma)$.
\item Step 3: $(E_0)$ homogeneous on $(R, a)$; b2b Steps 1--5 give
  $h = 0$ on $(R, a)$.
\item Step 4: $D(t) = 0$ for $t \in (0, R)$, and P1/P2 give
  $H = l = 0$ on $(0, R)$.
\item Step 5: reduce to A15.1b1 at $(R, T, T_0)$.
\end{itemize}
Hence $h = 0$.
\end{proof}

\begin{remark}[What remains open]
\label{rem:p12-b2d-remaining}
Corollary~\ref{cor:p12-b2d-wedge-plus} closes the region
$\sigma \le d/2$ of b2d completely.  The remaining open front
is
\[
\sigma > d/2,
\]
where core-both becomes nonempty: at $x \sim d/2$, both $H(x)$
and $H(d - x)$ are simultaneously live tail values.
b2d-upper is already treated by
Corollary~\ref{cor:p12-b2d-upper-local}, and b2d-core-single by
Theorem~\ref{thm:p12-b2d-core-single-full}.  So the residual obstruction
is precisely the symmetric core-both stratum.
\end{remark}

\section{A15.1b2d core-both closure and full Sub-wedge A}
\label{sec:p12-b2d-core-both}

We close the only remaining part of b2d.  Throughout this section let
\[
e/2 \le R < d/2,\qquad
R < \sigma < \varepsilon < \varepsilon_{\max},
\qquad
\sigma>d/2,
\]
and let \(h\in L^2(R,S)\) satisfy \(L_{T_0}h=0\), with
\(S=T+\sigma\), \(T_0=T+\varepsilon\).
Set
\[
I_{\rm both}:=
\bigl(\max\{R,d-\sigma\},\,\min\{\sigma,d-R\}\bigr).
\]
For \(x\in I_{\rm both}\) write
\[
y:=d-x,\qquad \kappa:=e-\delta.
\]
Then \(x,y\in(R,\sigma)\), and \(I_{\rm both}\) is invariant under
\(x\mapsto d-x\).

\begin{lemma}[Uniform core-both support bounds]
\label{lem:p12-b2d-core-both-support}
For a.e.\ \(x\in I_{\rm both}\),
\[
0<e-x<R,\qquad 0<e-y<R,
\]
\[
0<x-\delta<R,\qquad 0<y-\delta<R,
\]
and
\[
0<2\delta-x<R.
\]
Moreover
\[
y+\delta>\sigma,\qquad x+\kappa>\sigma.
\]
\end{lemma}

\begin{proof}
Since \(x,y\in(R,\sigma)\), \(R\ge e/2\), and
\(\sigma<\varepsilon_{\max}<e\),
\[
0<e-x<e-R\le e/2\le R,
\]
and the same holds with \(y\) in place of \(x\).

Next \(x,y>R\ge e/2>\delta\), while
\[
x-\delta,\ y-\delta
<\varepsilon_{\max}-\delta<e/2\le R.
\]
Also \(\varepsilon_{\max}<2\delta\), so \(2\delta-x>0\), and
\(R>\delta\) gives
\[
2\delta-x<2\delta-R<R.
\]

Finally
\[
y+\delta>e/2+\delta>\varepsilon_{\max}>\sigma.
\]
Since \(e>2\delta\), \(\kappa=e-\delta>\delta\), hence also
\[
x+\kappa>e/2+\delta>\varepsilon_{\max}>\sigma.
\]
The elementary comparisons used here are
\(e>2\delta\iff256>243\),
\(\varepsilon_{\max}<e\iff15<16\),
\(\varepsilon_{\max}<2\delta\iff80<81\), and
\(\varepsilon_{\max}-\delta<e/2\iff100<108\).
\end{proof}

\begin{theorem}[Full b2d core-both local kill]
\label{thm:p12-b2d-core-both-full}
Under the hypotheses above,
\[
h(x)=H(x)=l(x)=0
\]
for a.e.\ \(x\in I_{\rm both}\).
\end{theorem}

\begin{proof}
Each source map below is affine in \(x\) with derivative \(\pm1\).
Hence the preimage of a source-null set is null, and after intersecting
finitely many full-measure sets all source equations may be used
simultaneously for a.e.\ \(x\in I_{\rm both}\).

Consider the thirteen source positions
\[
\begin{aligned}
u_1&=d-x, &
u_2&=2d-x, &
u_3&=b-x, &
u_4&=3d-x,\\
u_5&=T-x, &
u_6&=T+\delta-x, &
u_7&=T+d-x,\\
u_8&=e+x, &
u_9&=a+x, &
u_{10}&=3e+x,\\
u_{11}&=b-\delta+x, &
u_{12}&=T-\delta+x, &
u_{13}&=T+x.
\end{aligned}
\]
They are all horizon-legal.  Indeed, the only sources above \(T\) are
\[
u_7=T+y<T+\sigma<T_0,\qquad
u_{12}=T+(x-\delta)<T+x<T+\sigma<T_0,
\]
and
\[
u_{13}=T+x<T+\sigma<T_0.
\]
All remaining sources are \(<T\): for example
\(u_2=d+y<d+e=a\),
\(u_4=2d+y<2d+e=b\),
\(u_6<T\) because \(x>\delta\),
\(u_9<a+e<T\), and
\(u_{10}<3e+\varepsilon_{\max}<4e<T\).

Applying
\[
Lh(u)=p[h(u-a)-h(u+a)]
+r[h(u-b)-h(u+b)]
+q[h(u-T)-h(u+T)]
\]
with the anti-reflection \(h(-t)=-h(t)\), zero extension outside
\((R,S)\), and Lemma~\ref{lem:p12-b2d-core-both-support}, gives:
\begin{align*}
Q_1:\;&-p h(e+x)-p h(b-x)-r h(a+x)-r h(T+\delta-x)\\
&\hspace{28mm}-q h(b-\delta+x)-q H(y)=0,\\
Q_2:\;&-p h(T+\delta-x)-r h(e+x)-q h(2e+x)=0,\\
Q_3:\;&p h(y)-pH(y)-r h(x)-q h(e+x)=0,\\
Q_4:\;&p h(y+\delta)-q h(x+\kappa)=0,\\
Q_5:\;&p h(a-x)-q h(x)=0,\\
Q_6:\;&p h(2d-x)+r h(y)=0,\\
Q_7:\;&p h(b-x)+r h(a-x)+q h(y)=0,\\
Q_8:\;&-p h(y)-p l(y)-r h(2d-x)-rH(x)-q h(b-x)=0,\\
Q_9:\;&p h(x)-pH(x)-r h(y)-q h(a-x)=0,\\
Q_{10}:\;&p h(x+\kappa)-q h(y+\delta)=0,\\
Q_{11}:\;&p h(e+x)-q h(y)=0,\\
Q_{12}:\;&p h(2e+x)+r h(x+\kappa)=0,\\
Q_{13}:\;&p h(a+x)+r h(e+x)+q h(x)=0.
\end{align*}
Here \(h(b-\delta+x)=h(T-y)=l(y)\).

Let \(\Delta:=p^2-q^2\).  From \(Q_4,Q_{10}\),
\[
\Delta\,h(y+\delta)=\Delta\,h(x+\kappa)=0.
\]
Since \(\Delta>0\),
\[
h(y+\delta)=h(x+\kappa)=0.
\]
Then \(Q_{12}\) gives \(h(2e+x)=0\).  The remaining equations yield
successively
\[
h(e+x)=\frac qp\,h(y),
\qquad
h(T+\delta-x)=-\frac{rq}{p^2}h(y),
\]
\[
h(2d-x)=-\frac rp\,h(y),
\qquad
h(a-x)=\frac qp\,h(x),
\]
\[
h(a+x)=-\frac qp\,h(x)-\frac{rq}{p^2}h(y),
\]
\[
h(b-x)=-\frac{rq}{p^2}h(x)-\frac qp\,h(y),
\]
\[
H(y)=\frac{\Delta}{p^2}h(y)-\frac rp\,h(x),
\qquad
H(x)=\frac{\Delta}{p^2}h(x)-\frac rp\,h(y).
\]
Substitution in \(Q_8\) gives
\[
l(y)=
-\frac{r(p^2-2q^2)}{p^3}h(x)
-\frac{\Delta-2r^2}{p^2}h(y).
\]
Finally \(Q_1\) reduces exactly to
\[
\frac{2qr(2p^2-q^2)}{p^3}\,h(x)=0.
\]
Since \(p,q,r>0\) and \(2p^2-q^2>0\),
\[
h(x)=0
\]
for a.e.\ \(x\in I_{\rm both}\).

The involution \(x\mapsto y=d-x\) preserves \(I_{\rm both}\) and
Lebesgue-null sets, so the same conclusion gives \(h(y)=0\).
Then \(Q_5,Q_9\) give \(h(a-x)=H(x)=0\), and the same argument at
\(y\) gives \(H(y)=0\).  Identity P1
(Theorem~\ref{thm:p12-P1}) therefore gives \(l(x)=0\).
\end{proof}

\begin{remark}[Polynomial row certificate]
\label{rem:p12-b2d-core-both-certificate}
The thirteen rows admit an independent exact certificate.  With
\(\Delta=p^2-q^2\), set
\[
\begin{aligned}
\widetilde C=(&
p^3\Delta,\,
-p^2r\Delta,\,
-p^2q\Delta,\,
pq^2r^2,\,
-r\Delta(p^2-2q^2),\\
&-pqr\Delta,\,
p\Delta^2,\,
-p^2q\Delta,\,
pqr\Delta,\,
p^2qr^2,\\
&p\Delta(\Delta-2r^2),\,
-pqr\Delta,\,
p^2r\Delta).
\end{aligned}
\]
Then, identically in all visibility variables,
\[
\sum_{j=1}^{13}\widetilde C_jQ_j
=
2qr\Delta(2p^2-q^2)\,h(x).
\]
Thus the triangular elimination and the row-multiplier calculation
cross-check the same non-vanishing closure.
\end{remark}

\begin{theorem}[Full A15.1b2d Sub-wedge A injectivity]
\label{thm:p12-b2d-full}
Let
\[
2a<T_0<c,\qquad e/2\le R<d/2,\qquad T<S<T_0.
\]
Then
\[
\ker L_{R,S,T_0}^{\{a,b,2a\}}=\{0\}.
\]
\end{theorem}

\begin{proof}
Write \(\sigma=S-T\).

If \(\sigma\le d/2\), this is
Corollary~\ref{cor:p12-b2d-wedge-plus}.  Assume therefore
\(\sigma>d/2\).

If \(d/2<\sigma\le d-R\), the defect strip decomposes, up to endpoints,
as
\[
(R,d-\sigma)\quad\text{(core-single)}
\]
and
\[
(d-\sigma,\sigma)\quad\text{(core-both)}.
\]
The first interval is killed by
Theorem~\ref{thm:p12-b2d-core-single-full}, including its visibility
values \(H,l\); the second is killed by
Theorem~\ref{thm:p12-b2d-core-both-full}.

If \(\sigma>d-R\), the defect strip decomposes as
\[
(R,d-R)\quad\text{(core-both)}
\]
and
\[
(d-R,\sigma)\quad\text{(b2d-upper)}.
\]
The first interval is killed by
Theorem~\ref{thm:p12-b2d-core-both-full}; the second by
Corollary~\ref{cor:p12-b2d-upper-local}.  Hence in either case
\[
h=H=l=0
\quad\text{a.e. on }(R,\sigma).
\tag{C}
\]

Consequently the \(E_\sigma\)-equation is homogeneous on all of
\((R,a)\).  The reflection Steps~1--5 of
Theorem~\ref{thm:p12-b2b} apply exactly as in the proof of
Corollary~\ref{cor:p12-b2d-wedge-plus}, and give
\[
h=0
\quad\text{a.e. on }(R,a).
\tag{E}
\]

It remains to kill the small tail.  Let \(0<t<R\).
In P2,
\[
D(t)=(p^2-q^2-r^2)h(t)+qr[h(e-t)-h(e+t)].
\]
Here \(h(t)=0\) by support, \(e+t\in(R,a)\) and is zero by (E), while
\(e-t\) is either below \(R\) or belongs to \((R,e)\subset(R,a)\).
Thus \(D(t)=0\), so P2 gives
\[
H(t)=l(t).
\tag{1}
\]

For P1 put \(z=d-t\).  Since
\[
z>d-R>R,
\]
either \(z\ge\sigma\), in which case \(H(z)=0\) by tail support, or
\(R<z<\sigma\), in which case \(H(z)=0\) by (C).  Hence always
\(H(d-t)=0\).  P1 gives
\[
H(t)+l(t)=0.
\tag{2}
\]
From (1)--(2),
\[
H(t)=l(t)=0
\quad\text{for a.e. }0<t<R.
\]
Together with (C), the entire tail \((T,S)\) vanishes.  The support
therefore reduces to \((R,T)\), and
Theorem~\ref{thm:p12-b1} at \((R,T,T_0)\) gives \(h=0\).
\end{proof}


\section{Consolidated status and the remaining wedge}\label{sec:p12-status}

\begin{corollary}[A15.1 consolidated injectivity]
\label{cor:p12-consolidated}
Let \(2a<T_0<c\) and \(0<R<S<T_0\).  Then
\[
\ker L_{R,S,T_0}^{\{a,b,2a\}}=\{0\}
\]
in the following established strata:
\begin{itemize}
\item[(i)] \(S<T\)
  (Theorem~\ref{thm:p12-b0});
\item[(ii)] \(S=T\)
  (Theorem~\ref{thm:p12-b1});
\item[(iii)] \(T<S<T_0\), \(e\le R<T\)
  (Theorem~\ref{thm:p12-b2a});
\item[(iv)] \(T<S<T_0\), \(d/2\le R<e\)
  (Theorem~\ref{thm:p12-b2c}, together with b2b where applicable);
\item[(v)] \(T<S<T_0\), \(e/2\le R<d/2\)
  (Theorem~\ref{thm:p12-b2d-full}).
\end{itemize}
Consequently,
\[
\boxed{
T<S<T_0,\qquad e/2\le R<T
\quad\Longrightarrow\quad
\ker L_{R,S,T_0}^{\{a,b,2a\}}=\{0\}.
}
\]
The paired-source identities P1 and P2
(Theorems~\ref{thm:p12-P1}, \ref{thm:p12-P2}) remain valid
unconditionally for every \(0<x<\varepsilon\).
\end{corollary}

\begin{remark}[b2d closure]
Theorem~\ref{thm:p12-b2d-core-both-full} closes the symmetric
core-both region that remained after the full core-single
Theorem~\ref{thm:p12-b2d-core-single-full} and the b2d-upper
Corollary~\ref{cor:p12-b2d-upper-local}.  Thus
Theorem~\ref{thm:p12-b2d-full} supersedes the earlier partial
b2d wedge statements as the final Sub-wedge~A result.
\end{remark}

\begin{openproblem}[Descent below \(R=e/2\)]
\label{open:p12-below-ehalf}
For \(2a<T_0<c\), \(0<R<e/2\), \(T<S<T_0\): decide whether
\[
\ker L_{R,S,T_0}^{\{a,b,2a\}}=\{0\}.
\]
\end{openproblem}

\section{Firewall}\label{sec:p12-firewall}

None of \S\ref{sec:p12-persistence}--\ref{sec:p12-status} makes a
polar-gauge, terminal-transport, Object-X or RH statement.  The R14
firewall from P11 remains unchanged: even a full closure of
Open Problem~\ref{open:p12-below-ehalf} would decide only the
localized-hub kernel triviality in the three-shift chamber; by R36-A1
(P11) this is equivalent to dense range of the annular hub projection
\(P_{\mathcal A}H_{T_0}\), a modulus-layer statement.  Neither
polar-gauge convergence nor strong terminal transport nor a canonical
mediator for Object~X follows.

\end{document}
